% Generated from the Lean dependency data; see the comment at the top of content.tex.
\chapter{Representations, characters and quasi-characters}
\label{chap:reps}
Background of Thesis Ch.~2--4: tempered representations, the multiplicity $m(\pi)$,
Harish-Chandra characters, quasi-characters and their germs.

\section{Representations and the multiplicity}
\begin{definition}[Unitary representations]
\label{UnitaryRep}
\lean{GSpinGGP.UnitaryRep}
\leanok
A strongly continuous unitary representation of a topological group on a complex Hilbert space
(Thesis \S3.1).
\end{definition}

\begin{definition}[Tempered representations]
\label{IsTempered}
\lean{GSpinGGP.IsTempered}
\leanok
\uses{UnitaryRep}
$\pi$ is irreducible and tempered: irreducible and weakly contained in the left regular representation
of $G$ (Thesis Ch.~1, \S3.1.2). The Thesis uses the growth condition on matrix coefficients and notes the
equivalence (Cowling--Haagerup--Howe); the weak-containment form is used here because it stays valid
when the centre of $G$ is noncompact.
\end{definition}

\begin{definition}[The multiplicity $m(\pi)$]
\label{mult}
\lean{GSpinGGP.mult}
\leanok
\uses{G, GGPData, H, UnitaryRep, xi}
For a unitary representation $\pi$ of $G(F)$, \[ m(\pi)=\dim\Hom_{H(F)}(\pi^\infty,\xi)\in\mathbb N\cup\{\infty\}, \]
where $\pi^\infty$ is the space of smooth vectors (Thesis \S5.5).
\end{definition}

\begin{literature}[Schur's lemma]
\label{LIT-dixmier_commutant_scalar}
\lean{GSpinGGP.dixmier_commutant_scalar}
\leanok
\uses{UnitaryRep}
For a continuous unitary representation of a topological group with no closed invariant subspaces other than $0$ and the
whole space, every bounded operator commuting with the group is a scalar.

\assumedfrom{\cite{Dixmier}, 13.1.4 (i)$\Rightarrow$(ii)}
\end{literature}

\begin{theorem}[The centre acts by a character]
\label{schur_central}
\lean{GSpinGGP.schur_central}
\leanok
\uses{UnitaryRep}
In an irreducible unitary representation of a topological group, the centre acts by a character (Thesis \S3.1.2).
\end{theorem}
\begin{proof}
\leanok
\uses{LIT-dixmier_commutant_scalar}
A central element's operator commutes with the whole group, so it is a scalar by Lemma~\ref{LIT-dixmier_commutant_scalar}.
\end{proof}

\begin{theorem}[$G(F)$ is unimodular]
\label{isMulRightInvariant_haar_G}
\lean{GSpinGGP.GGPData.isMulRightInvariant_haar_G}
\leanok
\uses{G, GGPData, GSpin}
The Haar measure of $G(F)=\GSpin(W)(F)\times\GSpin(V)(F)$ (and of every inner form) is right invariant.
\end{theorem}
\begin{proof}
\leanok
\uses{H, N, cartan_dieudonne_even}
The modular character $\Delta$ is a homomorphism to $\mathbb{R}_{>0}$, trivial on the centre and invariant under continuous
involutive automorphisms. Conjugation by an anisotropic vector $w$ on one factor is a continuous involutive automorphism sending $uw$ to
$q(u)q(w)(uw)^{-1}$, so $\Delta(uw)^2=1$ and $\Delta(uw)=1$. Scalars and products $uw$ of anisotropic vectors generate
$\GSpin(U)(F)$ (Cartan--Dieudonn\'e), so $\Delta=1$. (The thesis takes unimodularity as standard, \S2.1.)
\end{proof}

\begin{literature}[Weil's quotient-measure formula]
\label{LIT-weil_quotient}
\lean{GSpinGGP.weil_quotient}
\leanok
Let $S\hookrightarrow\Gamma$ be a closed embedding of unimodular locally compact, Hausdorff, second countable
groups with Haar measures $\mu,\nu$. There is a unique measure $m$ on $S\backslash\Gamma$ with
$\int_\Gamma\varphi\,d\nu=\int_{S\backslash\Gamma}\int_S\varphi(sq)\,d\mu(s)\,dm(q)$ for all
$\varphi\in C_c(\Gamma)$. (Used for $H(F)/A_{G,H}(F)\subset G(F)/A_G(F)$, Thesis \S7.1.)

\assumedfrom{Bourbaki, \emph{Int\'egration} VII \S2; \cite{Folland}, Thm.~2.49}
\end{literature}

\begin{literature}[Integration over $G=PK$]
\label{LIT-integral_eq_mul_integral_compact_integral}
\lean{GSpinGGP.integral_eq_mul_integral_compact_integral}
\leanok
Let $G$ be a locally compact, Hausdorff, second countable group, $P$ a closed subgroup and $K$ a compact
subgroup with $G=PK$. For Haar measures $dx,dp,dk$ there is $c>0$ with
$\int_G f(x)\,dx = c\int_K\int_P f(pk)\,dp\,dk$ for continuous compactly supported $f$.

\assumedfrom{\cite{Knapp}, Thm.~8.32; Bourbaki, \emph{Int\'egration} VII \S2; Casselman's notes \S3.1}
\end{literature}

\section{Characters and quasi-characters}
\begin{definition}[The Weyl discriminant]
\label{weylDisc}
\lean{GSpinGGP.weylDisc}
\leanok
$D^G(x)=|\det(1-\Ad x)|_{\mathfrak g/\mathfrak g_x}|_F$ (Thesis \S2.1), computed as the absolute value of the
trailing coefficient of the characteristic polynomial of $1-\Ad x$ on $\mathfrak g$.
\end{definition}

\begin{literature}[Harish-Chandra's regular set]
\label{LIT-hc_regular_open_dense_conull}
\lean{GSpinGGP.hc_regular_open_dense_conull}
\leanok
\uses{G, GGPData, GSpin, gspin}
For $x\in G(F)$ write $\det(t-1+\Ad(x))=D(x)t^\ell+\dots$ (terms of higher degree), $\ell=\operatorname{rank}G$. Then $G'=\{x: D(x)\neq0\}$ is open and dense in $G(F)$ and its complement has Haar measure zero.

\assumedfrom{\cite{HC70}, Part V \S3, p.~54 (also Introduction, p.~2), for $G$ connected reductive; here $G=\GSpin(W)\times\GSpin(V)$, $\ell=(\lfloor m/2\rfloor+1)+(\lfloor (m+2r+1)/2\rfloor+1)$}
\end{literature}

\begin{literature}[Centralizers in $\mathfrak{so}(V)$ have dimension at least the rank]
\label{LIT-jantzen_finrank_soCentMat_ge}
\lean{GSpinGGP.jantzen_finrank_soCentMat_ge}
\leanok
Over an algebraically closed field of characteristic $0$, for $n\ge3$ and $X\in\mathfrak{so}_n$, $\dim\mathfrak{so}_n(X)\ge\lfloor n/2\rfloor$.

\assumedfrom{\cite{Jantzen}, \S7.11, Remark, and \S7.13}
\end{literature}

\begin{literature}[Regular nilpotents of $\mathfrak{so}(V)$: Jordan type]
\label{LIT-jantzen_regularNil_partition}
\lean{GSpinGGP.jantzen_regularNil_partition}
\leanok
Over an algebraically closed field of characteristic $0$, for $n\ge3$, a nilpotent $X\in\mathfrak{so}_n$ whose centralizer has
minimal dimension among nilpotent elements has partition $[n]$ ($n$ odd) or $[n-1,1]$ ($n$ even): $X^{n-1}\neq0$, resp.\
$X^{n-2}\neq0=X^{n-1}$.

\assumedfrom{\cite{Jantzen}, \S6.3 (regular $\Leftrightarrow$ minimal centralizer dimension; partitions of the regular orbit)}
\end{literature}

\begin{lemma}[Centralizers of regular elements]
\label{finrank_lieCentralizer_eq_lieRank}
\lean{GSpinGGP.finrank_lieCentralizer_eq_lieRank}
\leanok
\uses{G, GGPData, gspin}
For $x\in\Greg(F)$, $\dim\mathfrak g_x=\operatorname{rank}\mathfrak g=\min_{X\in\mathfrak g}\dim\mathfrak g_X$.
\end{lemma}
\begin{proof}
\leanok
\uses{GSpin, LIT-jantzen_finrank_soCentMat_ge}
Both sides split over $\mathfrak g=\gspin(W)\times\gspin(V)$. For $g\in\GSpin(U)$, $\gspin(U)_g\cong F\oplus\mathfrak{so}(U)_{P(g)}$; if $g$ is semisimple with $\gspin(U)_g$ abelian then, over $\bar F$, the eigenvalues $\mu\ne\pm1$ of $P(g)$ are simple and pair off, $\dim E_{\pm1}\le2$, $\dim E_{-1}$ is even, and $\dim\mathfrak{so}(U)_{P(g)}=\lfloor\dim U/2\rfloor$ by induction over orthogonal splittings; the dimension descends to $F$. Finally $\operatorname{rank}\gspin(U)=\lfloor\dim U/2\rfloor+1$: Jantzen's lower bound \cite{Jantzen}, and an explicit skew element $\sum c_iJ_i$ on orthogonal planes.
\end{proof}

\begin{lemma}[$\Greg(F)\subseteq G'$]
\label{hcDisc_ne_zero_of_mem_regularSet}
\lean{GSpinGGP.hcDisc_ne_zero_of_mem_regularSet}
\leanok
\uses{G, GGPData, gspin}
For $x\in\Greg(F)$, $D(x)\neq0$.
\end{lemma}
\begin{proof}
\leanok
\uses{GSpin, finrank_lieCentralizer_eq_lieRank}
$N=1-\Ad x$ is semisimple on $\mathfrak g$, so $\mathfrak g=\ker N\oplus\operatorname{im}N$ with $\ker N=\mathfrak g_x$; the characteristic polynomial of $N$ is $t^{\dim\mathfrak g_x}$ times a polynomial with nonzero constant term $\pm\det(N|_{\operatorname{im}N})$, and $\dim\mathfrak g_x=\ell$.
\end{proof}

\begin{lemma}[$G'\subseteq\Greg(F)$]
\label{mem_regularSet_of_hcDisc_ne_zero}
\lean{GSpinGGP.mem_regularSet_of_hcDisc_ne_zero}
\leanok
\uses{G, GGPData, gspin}
If $D(x)\neq0$ then $x$ is semisimple with abelian centralizer in $\mathfrak g$.
\end{lemma}
\begin{proof}
\leanok
\uses{GSpin, LIT-jantzen_finrank_soCentMat_ge}
$D(x)\neq0$ says the generalized $1$-eigenspace of $\Ad x$ on $\mathfrak g$ has dimension $\le\ell$. For the Jordan decomposition $x=su$ in $G(F)$, $\mathfrak g_s$ lies in it, so $\dim\mathfrak g_s\le\ell$; the rank clauses for $\GSpin$ give $\dim\mathfrak g_s=\ell$, $\mathfrak g_s$ abelian, and the only unipotent element of $G(F)$ commuting with $s$ is $1$, so $u=1$ and $x=s$ is regular semisimple.
\end{proof}

\begin{theorem}[$G'=\Greg(F)$]
\label{hcDisc_ne_zero_iff_mem_regularSet}
\lean{GSpinGGP.hcDisc_ne_zero_iff_mem_regularSet}
\leanok
\uses{G, GGPData, gspin}
In characteristic $0$, $D(x)\neq0$ if and only if $x$ is semisimple with abelian centralizer in $\mathfrak g$.
\end{theorem}
\begin{proof}
\leanok
\uses{hcDisc_ne_zero_of_mem_regularSet, mem_regularSet_of_hcDisc_ne_zero}
The two inclusions above.
\end{proof}

\begin{theorem}[The regular set]
\label{hc_regularSet_open_conull}
\lean{GSpinGGP.hc_regularSet_open_conull}
\leanok
\uses{G, GGPData, GSpin, gspin}
The regular semisimple set $\Greg(F)$ is open in $G(F)$ and its complement has Haar measure zero.
\end{theorem}
\begin{proof}
\leanok
\uses{LIT-hc_regular_open_dense_conull, hcDisc_ne_zero_iff_mem_regularSet}
Harish-Chandra's theorem and the identification $G'=\Greg(F)$.
\end{proof}

\begin{literature}[Harish-Chandra's regularity theorem]
\label{LIT-hc_character_exists}
\lean{GSpinGGP.hc_character_exists}
\leanok
\uses{G, GGPData, GSpin, IsTempered, UnitaryRep, gspin}
For an irreducible tempered $\pi$, the distribution $f\mapsto\operatorname{tr}\pi(f)$ on $C_c^\infty(G(F))$
is given by a locally integrable function which is locally constant on $\Greg(F)$.

\assumedfrom{Bernstein (admissibility); \cite{HC99}; \cite{BP}, \S2.1; Thesis Ch.~1}
\end{literature}

\begin{definition}[The character $\theta_\pi$]
\label{charG}
\lean{GSpinGGP.charG}
\leanok
\uses{G, GGPData, GSpin, UnitaryRep, gspin, mult}
The Harish-Chandra character $\theta_\pi$ of $\pi$, the representative of $f\mapsto\operatorname{tr}\pi(f)$
that is locally constant on $\Greg(F)$ (Thesis Ch.~1).
\end{definition}

\begin{definition}[Fourier transforms of nilpotent orbits]
\label{jHat}
\lean{GSpinGGP.jHat}
\leanok
For a nilpotent orbit $\mathcal O\subset\mathfrak g(F)$, $\hat{j}(\mathcal O,\cdot)$ is the function
representing $f\mapsto\mu_{\mathcal O}(\hat f)$, the Fourier transform being normalized by $b$ and $\psi$
(\cite{BP} \S1.8, \S4.1; Thesis \S4.1.1). It is unique on $\mathfrak g_{\mathrm{reg}}$.
\end{definition}

\begin{definition}[Quasi-characters]
\label{IsQuasiCharacter}
\lean{GSpinGGP.IsQuasiCharacter}
\leanok
\uses{jHat}
$\theta: G(F)\to\mathbb C$ is a quasi-character (Thesis \S4.1.1, \cite{BP} \S4.1) if it is conjugation
invariant on $\Greg(F)$ and at every semisimple $x$ there are coefficients $c_{\theta,\mathcal O}(x)$,
$\mathcal O\in\Nil(\mathfrak g_x)$, with
$\theta(xe^X)=\sum_{\mathcal O}c_{\theta,\mathcal O}(x)\,\hat j(\mathcal O,X)$ for $X\in\mathfrak g_x(F)$
regular and near $0$.
\end{definition}

\begin{definition}[Germs]
\label{germ}
\lean{GSpinGGP.germ}
\leanok
\uses{jHat}
For a quasi-character $\theta$, semisimple $x$ and $\mathcal O\in\Nilreg(\mathfrak g_x)$, the germ
$c_{\theta,\mathcal O}(x)$ is the coefficient of $\hat j(\mathcal O,\cdot)$ in the local expansion of $\theta$
at $x$ (Thesis \S4.2, \cite{BP} \S4.5); it is $0$ if $\mathcal O$ is not regular or there is no expansion.
\end{definition}

\begin{literature}[Uniqueness of germs]
\label{LIT-germ_unique}
\lean{GSpinGGP.germ_unique}
\leanok
\uses{G, GGPData, gspin, jHat}
At semisimple $x$, the coefficients of a local expansion $\theta(xe^X)=\sum_{\mathcal O}c_{\mathcal O}\hat j(\mathcal O,X)$
are unique on the regular orbits $\mathcal O\in\Nilreg(\mathfrak g_x)$ (linear independence of the
$\hat j(\mathcal O,\cdot)$ near $0$).

\assumedfrom{\cite{HC99}; \cite{BP}, \S4.5; Thesis \S4.2}
\end{literature}

\begin{literature}[Local character expansion]
\label{LIT-hc_character_isQuasiCharacter}
\lean{GSpinGGP.hc_character_isQuasiCharacter}
\leanok
\uses{G, GGPData, GSpin, IsQuasiCharacter, IsTempered, UnitaryRep, gspin}
For irreducible tempered $\pi$, $\theta_\pi$ is a quasi-character.

\assumedfrom{\cite{HC99}; \cite{Wal10}, \S4; \cite{BP}, Prop.~4.1.1(i); Thesis Prop.~4.1.1.1(i)}
\end{literature}

\begin{literature}[Rodier's theorem]
\label{LIT-rodier_germ_one}
\lean{GSpinGGP.rodier_germ_one}
\leanok
\uses{G, GGPData, GSpin, IsQuasiSplit, IsTempered, UnitaryRep, WhittakerDatum, germ, gspin}
Let $G$ be quasi-split. For each $\mathcal O\in\Nilreg(\mathfrak g)$ there is a Whittaker datum $(U,\chi)$
such that for every irreducible tempered $\pi$, $c_{\theta_\pi,\mathcal O}(1)=1$ if $\pi$ has a
$(U,\chi)$-Whittaker model and $0$ otherwise (Thesis Prop.~4.4.1(i)).

\assumedfrom{\cite{Rodier}; \cite{MW87}; \cite{BP}, Prop.~4.8.1(i)}
\end{literature}

\begin{literature}[Self-dual Haar measures]
\label{LIT-exists_isSelfDual}
\lean{GSpinGGP.exists_isSelfDual}
\leanok
For a finite-dimensional $F$-vector space $V$ with a nondegenerate bilinear form $B$ there is a Haar measure
$\mu$ on $V$ for which Fourier inversion holds: $f(Z)=\int\hat f(Y)\psi(-B(Z,Y))\,d\mu(Y)$ with
$\hat f(Y)=\int f(X)\psi(B(X,Y))\,d\mu(X)$, for all $f\in C_c^\infty(V)$.

\assumedfrom{Tate's thesis; \cite{Weil}, Ch.~II \S5, Ch.~VII \S2; \cite{BP}, \S1.6; Thesis \S2.4}
\end{literature}

